% Detailed table following the compact summary in Appendix C.
% Summaries checked against the registered Lean statements and their supporting
% declarations; see TABLE4_LEAN_AUDIT.md for the row-by-row comparison.
\subsection{Hypotheses and conclusions}
\label{app:dk-result-summaries}

Table~\ref{tab:dk-result-summaries} summarizes the hypotheses and main
conclusions of the \DKResultCount{} results listed in
Table~\ref{tab:dk-full-inventory}, using the same IDs
\citep[Sections~2--8]{davis1970rotation}.
Hilbert spaces are separable and real or complex unless specified otherwise.
$\sym{NormN}$ denotes an arbitrary unitarily invariant norm on its domain of
finiteness; each displayed norm is evaluated on an operator in that domain.
$\|\cdot\|_{\mathrm{op}}$ and $\|\cdot\|_{\mathrm{HS}}$ denote the operator
and Hilbert--Schmidt norms. Self-adjoint operators may be unbounded:
operator identities hold on their domains, and residuals or right-hand sides
in norm estimates have bounded extensions.
For the ambient tangent bound, the tangent is defined as a bounded operator,
equivalently $\|\sym{Psource}-\sym{Qtarget}\|_{\mathrm{op}}<1$.

For the geometric results, $\sym{Psource},\sym{Qtarget}$ are orthogonal
projections, $\sym{SourceDirectRotation}$ is a direct rotation, and
$\sym{Wcompetitor}$ is any orthogonal or unitary transporter satisfying
$\sym{Wcompetitor}\sym{Psource}\sym{Wcompetitor}^*=\sym{Qtarget}$.
Write
\[
 D_0=\DKRange(\sym{Psource})\cap\DKKernel(\sym{Qtarget}),\qquad
 D_1=\DKKernel(\sym{Psource})\cap\DKRange(\sym{Qtarget}).
\]
The source's acute condition is $D_0=D_1=\{0\}$.
For geometric and displacement rows from \path{DK-3.3-prop} onward,
the ranges and kernels of
$\sym{Psource},\sym{Qtarget}$ have matching dimensions and
$\dim D_0=\dim D_1$; the displacement results also assume
the directed sine block
$(I-\sym{Qtarget})|_{\DKRange(\sym{Psource})}$ compact.
$\sym{Theta0},\Theta_1$ are the source and complementary angle operators,
$\sym{Theta}$ their ambient direct sum, and $\theta_j$ the ordered principal
angles. In the source decomposition the direct rotation has blocks
$\bigl(\begin{smallmatrix}C_0&-S_0^*\\S_0&C_1\end{smallmatrix}\bigr)$
with $C_j\ge0$; $J$ is its skew angle generator, so
$\sym{SourceDirectRotation}=e^{J\sym{Theta}}$.
For real unitaries, spectral half-plane conditions refer to their
complexifications.

For perturbation results, $\sym{A0}$ is the self-adjoint trial operator,
also the source block of $\sym{A}$ in the reducing decomposition;
$\sym{A1}$ is its complementary reference block when used.
The operators $\sym{Lambda0},\sym{Lambda1}$ are the exact blocks of
$\sym{Acurrent}=\sym{A}+\sym{H}$.
The maps $\sym{E0},\sym{F0}$ give source and target coordinates,
$\sym{H0},\sym{H1}$ are the diagonal perturbation blocks, and
$\sym{R}=\sym{Acurrent}\sym{E0}-\sym{E0}\sym{A0}$.
Coordinate maps are isometries except in the two lower-frame rows.
The exact frames have orthogonal ranges whose sum is the ambient space;
reducing projections preserve operator domains.
We abbreviate the two spectral hypotheses, for $\sym{delta}>0$, by
\begin{align*}
 \mathcal I_{\sym{delta}}(B,C):\quad&
 \sigma(B)\subset[\beta,\alpha],\quad
 \sigma(C)\cap(\beta-\sym{delta},\alpha+\sym{delta})=\varnothing,
 \quad\text{or the roles interchanged};\\
 \mathcal O_{\sym{delta}}(B,C):\quad&
 B\le\alpha I,\quad C\ge(\alpha+\sym{delta})I.
\end{align*}
The endpoints satisfy $\beta\le\alpha$ and are chosen for each row;
the spectral intervals may be
half-infinite.

\begingroup
\small
\setlength{\tabcolsep}{4pt}
\renewcommand{\arraystretch}{1.15}
\begin{longtable}{@{}>{\raggedright\arraybackslash}p{0.20\linewidth}>{\raggedright\arraybackslash}p{0.34\linewidth}>{\raggedright\arraybackslash}p{\dimexpr0.46\linewidth-4\tabcolsep\relax}@{}}
\caption{Hypotheses and main conclusions of the \DKResultCount{} Davis--Kahan results.}
\label{tab:dk-result-summaries}\\
\toprule
\textbf{ID} & \textbf{Hypotheses} & \textbf{Conclusion}\\
\midrule
\endfirsthead
\multicolumn{3}{@{}l}{\textit{Hypotheses and conclusions (continued)}}\\[2pt]
\toprule
\textbf{ID} & \textbf{Hypotheses} & \textbf{Conclusion}\\
\midrule
\endhead
\midrule
\multicolumn{3}{r@{}}{\textit{Continued on the next page}}\\
\endfoot
\bottomrule
\endlastfoot
\multicolumn{3}{@{}l}{\textbf{Headline perturbation bounds}}\\*\addlinespace[3pt]
\path{S2-sin-theta} &
 $\mathcal I_{\sym{delta}}(\sym{A0},\sym{Lambda1})$. &
 $\sym{delta}\sym{NormN}(\sin\sym{Theta0})\le\sym{NormN}(\sym{R})$.\\\addlinespace
\path{S2-tan-theta} &
 $\mathcal O_{\sym{delta}}(\sym{A0},\sym{Lambda1})$;
 Rayleigh--Ritz condition $\sym{H0}=0$. &
 $\sym{delta}\sym{NormN}(\tan\sym{Theta0})\le\sym{NormN}(\sym{R})$;
 $\sym{delta}\sym{NormN}(\tan\sym{Theta})\le\sym{NormN}(\sym{H})$.\\\addlinespace
\path{S2-sin-two-theta} &
 $\sym{A},\sym{Acurrent}$ have a common domain;
 $\sigma(\sym{Lambda0})\subset[\beta,\alpha]$;
 $\sigma(\sym{Lambda1})$ outside
 $(\beta-\sym{delta},\alpha+\sym{delta})$. &
 $\sym{delta}\sym{NormN}(\sin(2\sym{Theta0}))\le2\sym{NormN}(\sym{R})$;
 $\sym{delta}\sym{NormN}(\sin(2\sym{Theta}))\le2\sym{NormN}(\sym{H})$.\\\addlinespace
\path{S2-tan-two-theta} &
 $\mathcal O_{\sym{delta}}(\sym{A0},\sym{A1})$;
 bounded $\sym{H}$, $\sym{H0}=\sym{H1}=0$. &
 $\sym{delta}\sym{NormN}(\tan(2\sym{Theta0}))\le2\sym{NormN}(\sym{R})$;
 $\sym{delta}\sym{NormN}(\tan(2\sym{Theta}))\le2\sym{NormN}(\sym{H})$.\\\addlinespace[5pt]
\multicolumn{3}{@{}l}{\textbf{Two subspaces and direct rotations}}\\*\addlinespace[3pt]
\path{DK-3.1-prop} &
 $D_0=D_1=\{0\}$. &
 A unique direct rotation exists. Positive diagonal blocks characterize it
 among transporters.\\\addlinespace
\path{DK-3.2-prop} &
 General pair of projections. &
 A direct rotation exists iff $\dim D_0=\dim D_1$.
 On $D_0\oplus D_1$, $\sym{SourceDirectRotation}^2=-I$;
 uniqueness can fail.\\\addlinespace
\path{DK-3.3-prop} &
 A direct rotation $\sym{SourceDirectRotation}$; conversely a unitary
 square root with spectrum in the closed right half-plane that maps
 $D_0$ onto $D_1$. &
 $\sym{SourceDirectRotation}^2=(2\sym{Qtarget}-I)(2\sym{Psource}-I)$,
 and its spectrum lies in the closed right half-plane.
 Conversely, such a square root is a direct rotation.\\\addlinespace
\path{DK-3.4-prop} &
 $C_0^2\ge\tfrac12 I$;
 $Q_-=(2\sym{Psource}-I)\sym{Qtarget}(2\sym{Psource}-I)$. &
 $\sym{SourceDirectRotation}^2$ is a direct rotation from
 $\DKRange(Q_-)$ to $\DKRange(\sym{Qtarget})$.\\\addlinespace
\path{DK-3.1-thm} &
 Matching dimensions as above; angle operators
 $0\le\sym{Theta0},\Theta_1\le\pi/2$. &
 The four intersection dimensions and the generic angle's spectral
 multiplicity completely classify pairs up to isometric equivalence.
 Any angle data with matching nonzero multiplicities
 and the prescribed total dimension are realizable.\\\addlinespace
\path{DK-3.1-cor} &
 Classification hypotheses;
 $\sym{Psource}(I-\sym{Qtarget})\sym{Psource}$ compact. &
 Classification reduces to angle eigenvalues with multiplicity:
 $\theta_j\downarrow0$. The two angle operators may have different
 zero multiplicities.\\\addlinespace
\path{DK-3.5-prop} &
 A direct rotation and its angle operator $\sym{Theta}$; the maximal
 eigenspace characterization assumes the acute case. &
 $\sym{Theta}$ commutes with $\sym{Psource},\sym{Qtarget},J,
 \sym{SourceDirectRotation}$. If $x\ne0$ and
 $\sym{Theta}x=\theta x$, then
 $\angle(x,\sym{SourceDirectRotation}x)=\theta$.
 In the acute case these eigenspaces are the unique maximal reducing
 subspaces of the prescribed angle.\\\addlinespace
\path{DK-3.2-cor} &
 Interchange $\sym{Psource}$ and $\sym{Qtarget}$. &
 $\sym{Theta}$ is unchanged, $J$ changes to $-J$, and the corresponding
 direct rotation becomes $\sym{SourceDirectRotation}^*$.\\\addlinespace[5pt]
\multicolumn{3}{@{}l}{\textbf{Displacement extremality}}\\*\addlinespace[3pt]
\path{DK-4.1-prop} &
 Compact angle data; arbitrary transporter $\sym{Wcompetitor}$. &
 For the positive principal angles, there are orthonormal
 $v_j\in\DKRange(\sym{Psource})$ with
 $\angle(v_j,\sym{Wcompetitor}v_j)\ge\theta_j$.
 Each singular value of restricted displacement is minimized by
 $\sym{SourceDirectRotation}$, with value $2\sin(\theta_j/2)$.\\\addlinespace
\path{DK-4.1-cor} &
 Same displacement setting. &
 $\sym{NormN}((I-\sym{SourceDirectRotation})\sym{Psource})
 \le\sym{NormN}((I-\sym{Wcompetitor})\sym{Psource})$.\\\addlinespace
\path{DK-4.2-prop} &
 Any orthonormal basis $(v_j)$ of $\DKRange(\sym{Psource})$;
 arbitrary transporter $\sym{Wcompetitor}$. &
 $\sum_j\sin^2\angle(v_j,\sym{Wcompetitor}v_j)
 \ge\sum_j\sin^2\theta_j$, including an infinite right-hand side.\\\addlinespace
\path{DK-4.3-prop} &
 Same displacement setting; put
 $D(W)=(I-W)^*(I-W)$. &
 $\sym{NormN}(D(\sym{SourceDirectRotation}))
 \le\sym{NormN}(D(\sym{Wcompetitor}))$:
 the direct rotation minimizes squared full displacement.\\\addlinespace
\path{DK-4.4-prop} &
 Real space; largest principal angle at most $\pi/3$. &
 The printed assertion that $\sym{SourceDirectRotation}$ minimizes
 every $\sym{NormN}(I-\sym{Wcompetitor})$ is false.
 A real four-dimensional trace-norm counterexample is given in
 \cref{sec:prop44}; the Q-norm repair is stated there separately.\\\addlinespace[5pt]
\multicolumn{3}{@{}l}{\textbf{Sylvester estimates and spectral cutoffs}}\\*\addlinespace[3pt]
\path{DK-5.1-thm} &
 Banach spaces; bounded $C,X,Y$; $BX-XC=Y$;
 $\|C\|_{\mathrm{op}}\le a$,
 $\|B^{-1}\|_{\mathrm{op}}\le(a+\sym{delta})^{-1}$, $a\ge0$.
 The map norm is compatible with these operator norms. &
 $\sym{delta}\|X\|\le\|Y\|$.
 The coefficient hypotheses may be interchanged.
 An unbounded $B$ is closed and densely defined with bounded inverse.\\\addlinespace
\path{DK-5.2-thm} &
 Self-adjoint $B\ge(\gamma+\sym{delta})I$,
 $C\le\gamma I$; bounded $X,Y$;
 $X\DKDom(C)\subset\DKDom(B)$ and $BX-XC=Y$ there. &
 $\sym{delta}\sym{NormN}(X)\le\sym{NormN}(Y)$,
 including unbounded $B,C$.\\\addlinespace
\path{DK-5.1-lem} &
 Orthogonal projections $\Omega_t\to I$ strongly; bounded $K$. &
 $s_j(K\Omega_t)\to s_j(K)$ for each singular-value index $j$.\\\addlinespace[5pt]
\multicolumn{3}{@{}l}{\textbf{Generalized single-angle bounds}}\\*\addlinespace[3pt]
\path{DK-6.1-lem} &
 Each diagonal block of $K$ is dominated by the corresponding block of
 $L$ in every $\sym{NormN}$. &
 The same domination holds for their block-diagonal sums.
 The converse holds if the two $K$ blocks are equisingular and the two
 $L$ blocks are equisingular.\\\addlinespace
\path{DK-6.2-lem} &
 Orthogonal projections $\Omega,\Upsilon$; bounded $K$. &
 $\sym{NormN}(\Omega K\Upsilon+
 (I-\Omega)K(I-\Upsilon))\le\sym{NormN}(K)$.\\\addlinespace
\path{DK-6.1-prop} &
 $\sym{A},\sym{Acurrent}$ have a common domain and bounded difference
 $\sym{H}$;
 both $\mathcal I_{\sym{delta}}(\sym{A0},\sym{Lambda1})$ and
 $\mathcal I_{\sym{delta}}(\sym{A1},\sym{Lambda0})$. &
 $\sym{delta}\sym{NormN}(\sin\sym{Theta})\le\sym{NormN}(\sym{H})$.\\\addlinespace
\path{DK-6.1-thm} &
 $\sym{E0}^*\sym{E0}\ge\varepsilon^2I$, $\varepsilon>0$;
 $\mathcal I_{\sym{delta}}(\sym{A0},\sym{Lambda1})$.
 Source and target dimensions may differ. &
 $\sym{delta}\varepsilon\sym{NormN}(\sin\sym{Theta0})
 \le\sym{NormN}(\sym{R})$.
 The directed sine has the singular values of
 $\sym{Psource}(I-\sym{Qtarget})$.\\\addlinespace
\path{DK-6.2-thm} &
 Same lower-frame condition; only
 $\operatorname{dist}(\sigma(\sym{A0}),\sigma(\sym{Lambda1}))
 \ge\sym{delta}$. &
 $\sym{delta}\varepsilon\|\sin\sym{Theta0}\|_{\mathrm{HS}}
 \le\|\sym{R}\|_{\mathrm{HS}}$.
 If $\sym{R}$ has finite rank, also
 $\sym{delta}\varepsilon\|\sin\sym{Theta0}\|_{\mathrm{op}}
 \le\|\sym{R}\|_{\mathrm{op}}\sqrt{\operatorname{rank}\sym{R}}$.\\\addlinespace
\path{DK-6.3-thm} &
 Isometric Rayleigh--Ritz trial data;
 $\mathcal O_{\sym{delta}}(\sym{A0},\sym{Lambda1})$.
 The source dimension may be smaller than the target dimension. &
 $\sym{delta}\sym{NormN}(\tan\sym{Theta0})
 \le\sym{NormN}(\sym{R})$.\\\addlinespace
\path{DK-6.3-lem} &
 Bounded $K$; projections $\Gamma,\Psi$ of rank at most $n\ge1$;
 $K\Psi=\Gamma K\Psi$;
 $\sum_{j\le n}s_j(K\Psi)^2$\newline
 ${}>\sum_{j\le n}s_j(K)^2-\eta^2$, $\eta>0$. &
 $\|\Gamma K(I-\Psi)\|_{\mathrm{op}}<\eta$:
 a nearly maximizing finite-rank compression has small leakage.\\\addlinespace[5pt]
\multicolumn{3}{@{}l}{\textbf{Spectral selection and repulsion}}\\*\addlinespace[3pt]
\path{DK-8.1-thm} &
 Ordered reference gap
 $\mathcal O_{\sym{delta}}(\sym{A0},\sym{A1})$ with endpoint $\alpha$;
 bounded $\sym{H}$, $\sym{H0}=\sym{H1}=0$.
 For the equivalence, $\dim D_0=\dim D_1$. &
 $\sym{Theta}\le\pi/4$ iff
 $\sym{Lambda0}\le\alpha I$ and
 $\sym{Lambda1}\ge(\alpha+\sym{delta})I$; such a reducing target exists.
 With direct-rotation-aligned exact frames,
 $\sym{A1}-\alpha I\le C_1(\sym{Lambda1}-\alpha I)C_1$,
 as quadratic forms, with the lower-block analogue.
 In finite dimensions there are eigenvalue and symmetric-gauge
 repulsion inequalities.\\\addlinespace
\path{DK-8.2-thm} &
 Double-angle sine hypotheses; $\dim D_0=\dim D_1$;
 either bounded $\sym{H}$ with
 $\|\sym{H}\|_{\mathrm{op}}<\sym{delta}/2$, or bounded
 $\sym{R}$ with $\|\sym{R}\|_{\mathrm{op}}<\sym{delta}/2$;
 $\sigma(\sym{A0})\subset
 [\beta-\sym{delta}/2,\alpha+\sym{delta}/2]$. &
 The corresponding ambient or directed double-angle sine bound holds
 on the branch with angles strictly below $\pi/4$.\\
\end{longtable}
\endgroup
